\section{March 05}

    Recall: let \(\omega\) be a symplectic form on \(M\). 
    Since it is closed, it defines a de Rham cohomology class \([\omega]\in H^2_{dR}(M,\bbR)\).

    \begin{theorem}
        If \(M\) is compact with a symplectic form \(\omega\), then \([\omega]\neq 0\in H^2_{dR}(M,\bbR)\).
    \end{theorem}
    \begin{proof}
        Suppose the contrary that \([\omega]=0\). 
        Then there exists a 1-form \(\theta\) such that \(\omega=d\theta\). 
        Let \(n = 2m = \dim M\). 
        Consider the top form \(\omega^m = d(\theta\wedge \omega^{m-1})\).
        Integrating over \(M\), we get
        \[ \int_M \omega^m = \int_M d(\theta\wedge \omega^{m-1}) = \int_{\partial M} \theta\wedge \omega^{m-1} = 0, \]
        where we used Stokes' theorem and the fact that \(M\) is compact without boundary.
        However, since \(\omega\) is a symplectic form, \(\omega^m\) is a volume form, and thus \(\int_M \omega^m > 0\).
        This is a contradiction, hence \([\omega]\neq 0\).
    \end{proof}

    \begin{corollary}
        Let \(M\) be a compact symplectic manifold. 
        Then \(H^2_{dR}(M,\bbR)\neq 0\).
    \end{corollary}

    More generally, we have the following result.

    \begin{theorem}
        Let \(M\) be a compact symplectic manifold of dimension \(n = 2m\). 
        Then \(H^{2k}_{dR}(M,\bbR)\neq 0\) for all \(0 \leq k \leq m\).
    \end{theorem}
    \begin{proof}
        The case \(k=0\) is trivial.
        We show that \([\omega^k]\neq 0\in H^{2k}_{dR}(M,\bbR)\) for all \(0 < k \leq m\). 
        Suppose the contrary that \([\omega^k]=0\) for some \(k\). 
        Then there exists a \((2k-1)\)-form \(\theta\) such that \(\omega^k = d\theta\).
        Similar to the previous proof, we have
        \[ \int_M \omega^m = \int_M \omega^k \wedge \omega^{m-k} = \int_M d\theta \wedge \omega^{m-k} = \int_M d(\theta \wedge \omega^{m-k}) = \int_{\partial M} \theta \wedge \omega^{m-k} = 0, \]
        which is a contradiction.
    \end{proof}

    \begin{example}
        The \(n\)-dimensional sphere \(S^n\) admits a symplectic form if and only if \(n=2\). 
        In fact, if \(n=2\), then \(S^2\) is symplectic with the area form as the symplectic form. 
        If \(n\neq 2\), then \(H^2_{dR}(S^n,\bbR)=0\), so \(S^n\) cannot admit a symplectic form.
    \end{example}

    \begin{example}
        The manifold \(S^2 \times S^6\) does not admit a symplectic form. 
        This is because \(H^4_{dR}(S^2 \times S^6,\bbR) = 0\), so it has no symplectic form.
    \end{example}

    \begin{example}
        The manifold \(S^1 \times S^3\) does not admit a symplectic form since \(H^2_{dR}(S^1 \times S^3,\bbR) = 0\).
        It is interesting to note that \(S^1 \times S^3\) admits a complex structure.
        Indeed, we have 
        \[ S^3 \times S^1 \cong \left( \bbC^2 \setminus \{0\} \right) / \Gamma, \]
        where \(\Gamma\) is the group generated by the holomorphic map \((z_1,z_2) \mapsto (2z_1, 2z_2)\).
    \end{example}

    \begin{theorem}
        All orientable surfaces admit symplectic forms.
    \end{theorem}
    \begin{proof}
        Let \(M\) be an orientable surface. 
        Then \(M\) admits a Riemannian metric \(g\). 
        Let \(\omega\) be the area form of \(g\). 
        Then \(\omega\) is a symplectic form on \(M\).
    \end{proof}

    \begin{example}
        The vector space \(\bbR^{2n}\) admits a symplectic form given by
        \[ \omega = dx_1 \wedge dy_1 + \cdots + dx_n \wedge dy_n, \]
        where \((x_1,y_1,\ldots,x_n,y_n)\) are the standard coordinates on \(\bbR^{2n}\).
        This is called the \emph{standard symplectic form} on \(\bbR^{2n}\).
    \end{example}

    Recall that \(\bbC^n = \bbR^{2n}\) has a standard complex structure \(J\) and a standard Riemannian metric \(g\).
    One can check that the standard symplectic form \(\omega\) on \(\bbC^n\) is given by \(\omega(X,Y) = g(JX,Y)\) for all \(X,Y \in T_z\bbC^n\).
    Hence, we say that \((\bbC^n, \omega, J, g)\) is a Kähler manifold.

    \begin{example}
        Let \(M = \bbP^n(\bbC)\) be the complex projective space. 
        Recall that we have many descriptions of \(\bbP^n(\bbC)\):
        \begin{enumerate}
            \item \(\bbP^n(\bbC) = \{ \text{complex lines in } \bbC^{n+1} \}\);
            \item \(\bbP^n(\bbC) = \left(\bbC^{n+1} \setminus \{0\}\right) / \bbC^*\);
            \item \(\bbP^n(\bbC) = S^{2n+1} / S^1\).
        \end{enumerate}
        Recall that the \emph{Fubini-Study} metric on \(\bbP^n(\bbC)\) is given by the descending of the standard metric on \(S^{2n+1} \subset \bbC^{n+1}\) since the \(S^1\)-action on \(S^{2n+1}\) is isometric.
        More explicitly, let \(z \in S^{2n+1}\) with image \([z] \in \bbP^n(\bbC)\).
        Then the tangent space at \([z]\) is given by
        \[ T_{[z]}\bbP^n(\bbC) = \{ X \in \bbC^{n+1} : \langle X, z \rangle = \langle X, \im z \rangle = 0 \}, \]
        where \(\langle \cdot, \cdot \rangle\) is the standard Euclidean inner product on \(\bbC^{n+1} \cong \bbR^{2n+2}\).
        The Fubini-Study metric \(g\) on \(\bbP^n(\bbC)\) is given by \(g(X,Y) = \langle X, Y \rangle\) for all \(X,Y \in T_{[z]}\bbP^n(\bbC) \subset \bbC^{n+1}\).
        And the almost complex structure \(J\) on \(\bbP^n(\bbC)\) is also given by \(X \mapsto \im X\) for all \(X \in T_{[z]}\bbP^n(\bbC) \subset \bbC^{n+1}\).

        For every \(X,Y \in T_{[z]}\bbP^n(\bbC)\), we define the Fubini-Study symplectic form \(\omega\) on \(\bbP^n(\bbC)\) by
        \[ \omega(X,Y) = g(JX,Y) = \langle \im X, Y \rangle. \]
        This is a non-degenerate 2-form on \(\bbP^n(\bbC)\).
        \Yang{By direct computation}, we can check that \(\omega\) is closed, hence it is a symplectic form on \(\bbP^n(\bbC)\).
        In fact, \((\bbP^n(\bbC), \omega, J, g)\) is a Kähler manifold.
    \end{example}

    \begin{definition}
        Let \(M\) be a smooth manifold with a(n almost) complex structure \(J\), a Riemannian metric \(g\).
        We say that \((M, J, g)\) is an \emph{(almost) Kähler manifold} if the following conditions hold:
        \begin{enumerate}
            \item \(g(JX,JY) = g(X,Y)\) for all \(X,Y \in TM\);
            \item \(\omega(X,Y) \coloneqq g(JX,Y)\) is a closed 2-form on \(M\).
        \end{enumerate}
        The 2-form \(\omega\) is called the \emph{(almost) Kähler form} of \((M, J, g)\).
        It is easy to check that \(\omega\) is a symplectic form on \(M\).
    \end{definition}

    Assume that \((M, J, g)\) is a(n almost) K\"ahler manifold. 
    Let \(N\) is a(n almost) complex submanifold of \(M\). 
    Here we say that \(N\) is a(n almost) complex submanifold of \(M\) if \(J(TN) \subset TN\).
    Then the restriction of \(\omega\) to \(N\) is a(n almost) K\"ahler form on \(N\).
    Note that the submanifold \(N\) of a symplectic manifold \((M, \omega)\) is not necessarily a symplectic submanifold of \(M\) 
    since the restriction of \(\omega\) to \(N\) may be degenerate.

    \begin{example}
        All complex submanifolds of \(\bbC^n\) and \(\bbP^n(\bbC)\) are K\"ahler manifolds and hence symplectic manifolds.
    \end{example}

    \begin{example}
        Let \(M\) be an arbitrary smooth manifold of dimension \(m\).
        Then the cotangent bundle \(T^*M \eqqcolon N\) admits a natural symplectic form \(\omega\).
        Let \(\pi: N \to M\) be the projection map.
        For every \(p \in M\) and \(\varphi \in T^*_pM\), and for \(v \in T_\varphi N\), \(\pi_*(v) \in T_pM\) is well-defined.
        We define 
        \[ \theta(v\in T_\varphi N) = \varphi(\pi_*(v)) \in \bbR. \]
        Then \(\theta\) is a 1-form on \(N\), and we define \(\omega = d\theta\).
        Easily see that \(\omega\) is closed.

        For the non-degeneracy of \(\omega\), we can work in local coordinates.
        Let \((U, x_1, \ldots, x_m)\) be a local coordinate chart on \(M\). 
        % Then we have a local coordinate chart \((\pi^{-1}(U), x_1, \ldots, x_m, y_1, \ldots, y_m)\) on \(N\), 
        % where \(y_i\) is the fiber coordinate corresponding to \(x_i\).
        For every \(\varphi \in T*U \subset T*M\), we can write \(\varphi = \sum_{i=1}^m y_i dx_i\) for some unique tuple \(y_1, \ldots, y_m \in \bbR\).
        Then we have a local coordinate chart \((\pi^{-1}(U), x_1, \ldots, x_m, y_1, \ldots, y_m)\) on \(N\).
        In this local coordinate chart, we have \(\theta = \sum_{i=1}^m y_i dx_i \in T^*N\).
        Hence, we have
        \[ \omega = d\theta = \sum_{i=1}^m dy_i \wedge dx_i. \]
        This is a non-degenerate 2-form on \(N\).
    \end{example}

    \begin{example}
        If \(M_i, \omega_i\) are symplectic manifolds for \(i=1,2\), then the product manifold \(M_1 \times M_2\) admits a natural symplectic form \(\omega = \pi_1^*\omega_1 + \pi_2^*\omega_2\), 
        where \(\pi_i: M_1 \times M_2 \to M_i\) is the projection map for \(i=1,2\).
    \end{example}




